\section{Qualitative Reconstruction Comparison}

\subsection{Original qualitative evidence}
\begin{draftpoints}
  \item original paper presents reconstructed digits after each new class and reports that NDL+IR retains recognizable previously learned digits more consistently than CL+IR \parencite{draelos2017neurogenesis}
  \item These panels are important because a small mean error difference may correspond either to visibly improved strokes or to diffuse pixel-level changes that are difficult to interpret semantically
  \item original images do not provide a controlled local comparison because they use the authors' models, samples, and undocumented training choices
\end{draftpoints}

\begin{figure}[htbp]
\centering
\includegraphics[width=0.95\textwidth]{original_reconstructions.png}
\caption{Original paper reconstruction comparison. Panel A shows CL+IR and panel B NDL+IR as classes are introduced; this figure is reproduced as the qualitative target, not as local replication evidence \parencite{draelos2017neurogenesis}.}
\label{fig:original-reconstructions}
\end{figure}

\subsection{Preserved local diagnostic}
\begin{draftpoints}
  \item workspace preserves one local MLflow partial-reconstruction diagnostic, which confirms that the logging pipeline can hold fixed samples and reconstruct them through multiple depths
  \item diagnostic is not a matched six-condition MNIST result and must not be used to rank CL, NDL, or replay regimes
  \item noisy lower-depth outputs illustrate why qualitative panels need depth labels, fixed sample identities, and per-image error values rather than unannotated example images
\end{draftpoints}

\begin{figure}[htbp]
\centering
\includegraphics[width=0.82\textwidth]{local_partial_reconstruction_diagnostic.png}
\caption{Only preserved local partial-reconstruction diagnostic. Columns are divided into learned and unlearned samples and rows show partial reconstructions. This artifact demonstrates instrumentation only; it is not a confirmatory comparison between methods.}
\label{fig:local-qualitative-diagnostic}
\end{figure}

\subsection{Required matched reconstruction grid}
\begin{draftpoints}
  \item definitive qualitative comparison requires rerunning all primary conditions with the same fixed test indices, preprocessing, seed policy, and curriculum checkpoints
  \item Grid rows should contain the input, base model, CL no replay, NDL no replay, CL original-data replay, NDL original-data replay, CL intrinsic replay, NDL intrinsic replay, and optimized NDL
  \item Columns should contain fixed examples from base classes, early incoming classes, late incoming classes, and deliberately difficult high-error examples
  \item Each panel should record condition, seed, class, sample index, curriculum step, depth, full-image MSE, and foreground MSE
  \item \missingartifact{The required matched checkpoint set and reconstruction PNGs are not preserved in the repository; the protocol therefore specifies the panel but does not fabricate it from unrelated runs.}
\end{draftpoints}

\begin{figure}[htbp]
\centering
\begin{tikzpicture}[
cell/.style={draw,minimum width=16mm,minimum height=8mm,font=\tiny,align=center},
label/.style={font=\tiny,align=right},
missing/.style={draw,fill=ProtocolLightPurple,minimum width=16mm,minimum height=8mm,font=\tiny,align=center}
]
\foreach \x/\name in {0/Base 1,1/Base 7,2/Early 0,3/Middle 4,4/Late 9,5/Hard case}{
  \node[cell] at (1.6*\x,0) {\name};
}
\foreach \y/\name in {1/Input,2/Base model,3/CL no replay,4/NDL no replay,5/CL original,6/NDL original,7/CL intrinsic,8/NDL intrinsic,9/Optimized NDL}{
  \node[label,anchor=east] at (-1.0,-0.8*\y) {\name};
  \foreach \x in {0,...,5}{\node[missing] at (1.6*\x,-0.8*\y) {fixed\\sample};}
}
\end{tikzpicture}
\caption{Planned matched reconstruction grid. Every method must reconstruct the same fixed images at the same curriculum checkpoint; cells remain a documented rerun requirement rather than being filled with unmatched examples.}
\label{fig:planned-reconstruction-grid}
\end{figure}

\begin{table}[htbp]
\centering
\caption{Manifest required for every qualitative reconstruction panel.}
\label{tab:qualitative-manifest}
\scriptsize
\begin{tabularx}{\textwidth}{p{2.8cm}Y}
\toprule
Field group & Required fields \\
\midrule
Panel identity & figure, row, column, condition, seed \\
Sample identity & dataset split, class ID, immutable sample index, learned/unlearned status \\
Training state & curriculum step, most recently learned class, architecture widths \\
Reconstruction & network depth, full-image MSE, foreground MSE, input path, output path \\
Provenance & checkpoint path, configuration snapshot, code revision \\
\bottomrule
\end{tabularx}
\end{table}

\subsection{Qualitative interpretation rules}
\begin{draftpoints}
  \item Visual comparisons must be made from fixed samples selected before inspecting model outputs; cherry-picked best or worst examples should be reported separately
  \item Recognition-like judgments should remain descriptive because the replication does not include a blinded human-rating protocol or a fixed downstream classifier
  \item Qualitative claims should agree with per-image and class-level error values; disagreements should be discussed as evidence that pixel MSE and perceived stroke quality capture different properties
\end{draftpoints}
